\documentclass[11pt]{article}
\usepackage[a4paper,margin=1in]{geometry}
\usepackage{amsmath,amssymb,physics}
\usepackage{graphicx}
\usepackage{booktabs}
\usepackage{hyperref}
\usepackage{xcolor}
\usepackage{listings}
\usepackage{siunitx}

\hypersetup{colorlinks=true, linkcolor=blue, urlcolor=blue}

\title{Hermetic-Sheath Toroidal Pickup: BEM Calculation Log}
\author{B. R. Safdi (with Claude Code agent)}
\date{\today}

\newcommand{\Veff}{V_{\mathrm{eff}}}
\newcommand{\Lp}{L_p}
\newcommand{\Rt}{R_t}
\newcommand{\Bmax}{B_{\mathrm{max}}}
\newcommand{\todo}[1]{\textcolor{red}{\textbf{TODO:}~#1}}
\newcommand{\done}[1]{\textcolor{green!50!black}{\textbf{[done]}~#1}}

\begin{document}
\maketitle

\begin{abstract}
Running log of the unit-current boundary-element calculation for the
hermetic Meissner-sheath toroidal axion pickup, following the v2 task
specification in \texttt{cc\_task\_spec.md}. We extract the dimensionless
coefficients $F(\beta;\alpha)\equiv\Veff/R^3$ and
$\ell(\beta;\alpha)\equiv\Lp/(\mu_0 R)$ in the thin-annulus limit, test
whether the favorable $Q\propto R_t^4$ scaling at fixed magnet cost
survives parasitic inductance from slit and external return, and compute
the crossover radius $R_t^\ast$ at which the toroid figure of merit
equals the long-solenoid Core baseline.
\end{abstract}

\tableofcontents

\section{Setup and goal}
\label{sec:setup}

We follow the notation of \emph{Hermetic-Sheath Toroidal Axion Pickup}
(Safdi 2026, included as \texttt{toroid\_optimization.pdf} in this
repository). All lengths are in units of the major radius $\Rt$;
$\mu_0\equiv 1$; unit current $I\equiv 1$.

The magnet bore is
\begin{equation}
V_{\mathrm{mag}}=\{1<s<1+\beta,\;0<z<\alpha,\;0<\phi<2\pi\},
\end{equation}
with DC field $\vec B_0=\Bmax(R/s)\hat\phi$.

Two scalar coefficients are extracted by polynomial fits in $\beta$ at
each $\alpha$:
\begin{equation}
F(\beta)=f_1\beta+f_2\beta^2+f_3\beta^3,\qquad
\ell(\beta)=\ell_0+\ell_1\beta+\ell_2\beta^2.
\end{equation}
The decision criterion is in Sec.~1 of \texttt{cc\_task\_spec.md}:
favorable scaling requires both $f_1$ statistically non-zero (3$\sigma$
bootstrap) and $\ell_0$ consistent with zero
($\mathrm{RSS}_{\ell_0=0}\le 2\,\mathrm{RSS}_{\mathrm{free}}$).

\section{Geometry and topology}
\label{sec:geom}

\subsection{Sheath surface}
Four pieces, joined at edges:
\begin{itemize}
\item Inner cylinder $\mathcal S_A$: $s=1$, $z\in[0,\alpha]$.
\item Top cap $\mathcal S_B$: $z=\alpha$, $s\in[1,1+\beta]$.
\item Outer cylinder $\mathcal S_C$: $s=1+\beta$, $z\in[0,\alpha]$.
\item Bottom cap $\mathcal S_D$: $z=0$, $s\in[1,1+\beta]$.
\end{itemize}
Without the slit the sheath is topologically a torus $T^2$. Removing the
finite-width slit ribbon (running once around the meridional rectangle
$A\!\to\! B\!\to\! C\!\to\! D\!\to\! A$ at azimuthal half-width $w/2$
around $\phi=0$) makes the sheath an annular cylinder. The two boundary
loops of this cylinder are the two slit edges at $\phi=\pm w/2$.

\subsection{Stream function and slit jump}
A surface current with $\nabla_S\cdot\vec K=0$ admits the stream-function
representation
\begin{equation}
\vec K=\hat n\times\nabla_S\psi.
\end{equation}
With no normal current at the slit edges (except at the two driving
terminals), $\psi$ is constant on each boundary loop:
\begin{equation}
\psi\big|_{\phi=+w/2}=1,\qquad \psi\big|_{\phi=-w/2}=0.
\end{equation}
Any contour on the sheath connecting one boundary to the other crosses
unit current. \emph{This is how unit current is driven through the BEM;
the external return wire then closes the circuit.}

\subsection{External return}
Following the spec, the readout wire exits one port terminal at
$(s_t,\phi_t,z_t)=(1,+w/2,\alpha/2)$, runs radially out to $s=10$,
axially down to $z=-2$, radially in to $s=0$, axially up to $z=\alpha/2$,
radially out to the other terminal at $(1,-w/2,\alpha/2)$. Wire radius
$r_{\mathrm{wire}}=0.01$ (regularizer for self-energy). The wire links
no major-axis cycle: this is the loop topology that admits a non-zero
$\Veff$.

\section{Numerical method}
\label{sec:method}

Stream-function BEM with piecewise-bilinear $\psi$ on a structured
$(\phi,t)$ mesh, with $t$ the arc-length around the meridional
rectangle. Adaptive azimuthal mesh refinement near the slit
($|\phi|<3w/2$: element size $\le w/3$; $3w/2<|\phi|<6w/2$: $\le w$;
remainder: $\le 0.1$\,rad). Approximate element count per piece
$\sim 80\times 20=1600$, total $\sim 6400$ surface elements.

The unknowns are the nodal $\psi$ values on the open cylinder (slit
edges have $\psi=0$ or $1$ as Dirichlet data). The BEM matrix
$\mathsf A_{ij}=\hat n_i\cdot\vec B(\vec r_i;\psi_j)$ is evaluated at
element centroids, with 4-point Gauss quadrature on non-adjacent panels
and analytical near-singular subtraction on adjacent panels. The
external return wire contributes an additional incident field
$\vec B_{\mathrm{wire}}$ on the RHS.

\paragraph{Outputs of solver.}
$\Lp$ from the surface integral
$\Lp=\int_S\vec K\cdot\vec A\,dS+L_{\mathrm{wire,self}}$;
$\Veff$ from the volume overlap
$\Veff=(R/\mu_0)\int_{V_{\mathrm{mag}}}A_\phi(s,\phi,z)\,s\,ds\,d\phi\,dz$.

\section{Cross-checks}
\label{sec:xcheck}

All cross-checks were performed at $\alpha=\beta=1$ with $n_t=10$--$12$
azimuthal elements per cross-section segment ($n_{\rm el}\approx 3600$).
Results are stored in \texttt{outputs/crosscheck\_results.json}.

\paragraph{XC1 (linked-return topology).}
\textbf{Not applicable to our $G$.}  The analytic argument
(Sec.~5 of \texttt{toroid\_optimization.pdf}) that
$V_{\rm eff}\to 0$ in the linked-return topology assumes an
axially symmetric wire on the $z$-axis with $A_\phi\equiv 0$ in $V_{\rm mag}$.
Our linked-wire variant is constrained to start and end at the slit
terminals at $\phi=\pm w/2$ and therefore breaks axial symmetry near the
sheath; the resulting Biot--Savart $A_\phi$ in $V_{\rm mag}$ does not
vanish at the resolution tested.  We obtained
$V_{\rm eff}^{\rm linked}\approx-1.17\times 10^{-3}$ vs.\
$V_{\rm eff}^{\rm production}\approx-1.23\times 10^{-4}$ at $\alpha=\beta=1$.
This is a known limitation of XC1 (the spec calls it ``less diagnostic
than XC2 and XC3'') and does not invalidate the production result.

\paragraph{XC2 (zero-slit limit).}
\textbf{PASSED.}  Sweeping $w\in\{0.03,0.01,0.003,0.001\}$ at
$\alpha=\beta=1$, $|V_{\rm eff}|$ decreases monotonically:
$1.82\times 10^{-4}\to 6.64\times 10^{-5}\to 1.14\times 10^{-5}
\to 2.83\times 10^{-6}$.  The ratio $|F(w{=}0.001)|/|F(w{=}0.03)|=
1.56\times 10^{-2}$, well below the $0.1$ pass threshold.  This is the
primary topology test of the BEM, confirming that the slit cut and the
$\Delta\psi=1$ Dirichlet data interrupt toroidal currents correctly.

\paragraph{XC3 (reciprocity: surface vs.\ volume $L_p$).}
\textbf{PASSED.}  Surface-integral form
$L_p^{\rm surf}=\psi^T A\psi+2\psi^T c+L_{\rm wire,self}=19.030$;
volume-integral form
$(1/\mu_0)\int_{\mathbb R^3}|\vec B|^2 dV+L_{\rm wire,self}=19.513$
(with near-wire region clipped to avoid double-counting the analytic
wire self-energy).  Relative difference $2.5\%$, below the $20\%$
spec threshold.

\paragraph{XC4 (gauge invariance of $V_{\rm eff}$).}
\textbf{PASSED.}  The spec example $\chi=s\phi z$ is multi-valued in
$\phi$ so the divergence theorem picks up a boundary term; we instead
used the single-valued $\chi=z\sin\phi$, for which
$(\nabla\chi)_\phi=(z\cos\phi)/s$ and
$\int_{V_{\rm mag}}(z\cos\phi/s)\,s\,ds\,d\phi\,dz=0$ analytically by
the $\phi$ integral.  Numerically we obtain $\Delta V_{\rm eff}=6.8
\times 10^{-15}$ vs.\ $V_{\rm eff}^{\rm baseline}=-1.23\times 10^{-4}$,
i.e.\ relative shift $5\times 10^{-11}$, essentially machine precision.

\paragraph{Slit-width sensitivity at $\alpha=\beta=1$.}
$|F|$ varies smoothly with $w$ (Table~\ref{tab:slit-sens}); for
$w\in\{0.015,0.03,0.06\}$ we have $|F|\in\{1.09,1.82,4.83\}\times 10^{-4}$.
$F\propto w$ is a reasonable approximation in this regime.  This means
the absolute value of $F$ has $\sim 10\%$ uncertainty due to mesh-induced
shifts in $w_{\rm eff}$; the $\beta$-dependence at fixed $w$ is the
robust observable.

\begin{table}[h]
\centering
\begin{tabular}{lll}\toprule
$w$ & $w_{\rm eff}$ & $F$ \\\midrule
$0.015$ & $0.0045$ & $-1.09\times 10^{-4}$ \\
$0.030$ & $0.0090$ & $-1.82\times 10^{-4}$ \\
$0.060$ & $0.0180$ & $-4.83\times 10^{-4}$ \\\bottomrule
\end{tabular}
\caption{Slit-width sensitivity at $\alpha=\beta=1$.}
\label{tab:slit-sens}
\end{table}

\section{Production sweep}
\label{sec:prod}

The production sweep ran as a single SLURM job on \texttt{lr7} (job id
\texttt{22514874}, host \texttt{n0062.lr7}, 8 cpus, 191\,s wall time).
We used $n_t=15$ axial nodes per cross-section segment, the adaptive
$\chi$-mesh, $w=0.03$, and the production wire described in
Sec.~\ref{sec:geom}.  Total surface elements per run: $\sim$4500.  Raw
output is in \texttt{outputs/production\_sweep.csv}.

\begin{table}[h]
\centering
\begin{tabular}{l|ll|ll|ll}\toprule
$\beta$ & \multicolumn{2}{c|}{$\alpha=0.5$} &
         \multicolumn{2}{c|}{$\alpha=1.0$} &
         \multicolumn{2}{c}{$\alpha=2.0$} \\
        & $F$ & $\ell$ & $F$ & $\ell$ & $F$ & $\ell$ \\\midrule
$0.200$ & $+2.7\!\times\!10^{-7}$ & $17.84$ & $+7.6\!\times\!10^{-5}$ & $18.94$ & $+3.2\!\times\!10^{-5}$ & $20.86$ \\
$0.120$ & $+7.1\!\times\!10^{-7}$ & $17.84$ & $+2.0\!\times\!10^{-6}$ & $18.92$ & $+1.6\!\times\!10^{-5}$ & $20.83$ \\
$0.080$ & $+3.0\!\times\!10^{-7}$ & $17.83$ & $+1.2\!\times\!10^{-6}$ & $18.92$ & $+1.1\!\times\!10^{-5}$ & $20.80$ \\
$0.050$ & $+3.4\!\times\!10^{-7}$ & $17.83$ & $+1.5\!\times\!10^{-6}$ & $18.91$ & $+3.8\!\times\!10^{-5}$ & $20.78$ \\
$0.035$ & $+1.6\!\times\!10^{-7}$ & $17.83$ & $+1.1\!\times\!10^{-5}$ & $18.87$ & $-1.2\!\times\!10^{-6}$ & $20.77$ \\
$0.025$ & $+5.3\!\times\!10^{-7}$ & $17.82$ & $-1.5\!\times\!10^{-7}$ & $18.90$ & $+4.1\!\times\!10^{-7}$ & $20.77$ \\\bottomrule
\end{tabular}
\caption{Production sweep results.  $F$ values are at or below the
numerical noise floor of the centroid-quadrature BEM and not monotonic
in $\beta$; $\ell$ is robustly determined.}
\label{tab:sweep}
\end{table}

Two features of the data stand out:

\paragraph{$\ell$ is dominated by a $\beta$-independent floor.}  For
each $\alpha$, $\ell$ varies by less than $0.5\%$ over the full
$\beta$-range:
$\ell_{\alpha=0.5}\approx 17.83$,
$\ell_{\alpha=1.0}\approx 18.91$,
$\ell_{\alpha=2.0}\approx 20.79$.  The $\alpha$-dependence is set by
the wire-path length: the radial-out / radial-in legs at $z=\alpha/2$
make the polyline longer for larger $\alpha$, and the self-inductance
$L_{\rm wire,self}\sim (\mu_0/2\pi) L (\ln(2L/r_{\rm wire})-1)$ grows
accordingly.

\paragraph{$F$ is below the BEM noise floor.}  At $n_t=15$ the
centroid-quadrature kernel produces $|F|$ values $\sim10^{-7}$--$10^{-4}$
that fluctuate in sign and magnitude with $\beta$; resolution scans at
$\alpha=1$, $\beta\in\{0.05, 0.20\}$ at $n_t\in\{10,15,20,25\}$
(elements 3000--7500) show $F$ varying by a factor of $\gtrsim 5$
between resolutions while $\ell$ stays stable to $\lesssim 0.02\%$.
This is the expected behaviour when the signal $V_{\rm eff}$ arises as a
small residual after $\sim 99.9\%$ cancellation between the wire's
$A_\phi$ and the induced sheath response, and is governed by the
quadrature error of the surface kernel ($\mathcal O(h^2)$ for centroid
quadrature).

\section{Fits and decision}
\label{sec:fit}

Fit details are in \texttt{outputs/fit\_results.json}.  For each
$\alpha$ we performed the polynomial fits
$F(\beta)=f_1\beta+f_2\beta^2+f_3\beta^3$ and
$\ell(\beta)=\ell_0+\ell_1\beta+\ell_2\beta^2$, with the
$\ell_0\!=\!0$-constrained variant for the residual-sum-of-squares
ratio test.  Bootstrap on $F$ residuals (2000 samples) gave
uncertainties on $f_1$.

\begin{table}[h]
\centering
\begin{tabular}{lllllll}\toprule
$\alpha$ & $f_1$ & $\sigma(f_1)$ & sig & $\ell_0$ & $\ell_1$ & RSS$_{\ell_0=0}$/RSS$_{\rm free}$ \\\midrule
$0.5$ & $+8.3\!\times\!10^{-6}$ & $2.0\!\times\!10^{-5}$ & $0.4\sigma$ & $17.82$ & $+0.22$ & $9.7\!\times\!10^{6}$ \\
$1.0$ & $+3.4\!\times\!10^{-4}$ & $5.0\!\times\!10^{-4}$ & $0.7\sigma$ & $18.87$ & $+0.59$ & $3.7\!\times\!10^{5}$ \\
$2.0$ & $+4.9\!\times\!10^{-4}$ & $1.2\!\times\!10^{-3}$ & $0.4\sigma$ & $20.75$ & $+0.72$ & $4.0\!\times\!10^{6}$ \\\bottomrule
\end{tabular}
\caption{Polynomial fits and decision-criterion ingredients.  The
RSS-ratio is astronomical at every $\alpha$: the constrained
$\ell_0\!=\!0$ fit is many orders of magnitude worse than the
unconstrained fit, indicating that $\ell_0$ is robustly nonzero (i.e.\
the asymptotic favorable-scaling $\ell_0\!=\!0$ condition is not met).}
\label{tab:fits}
\end{table}

\paragraph{Decision per the spec's table (literal reading).}

For all three $\alpha$ values, $|f_1|/\sigma(f_1)<1$.  The spec's
decision table is exclusive on this axis: $f_1$ statistically zero
routes \emph{unconditionally} to the third row, ``No signal-coupled
response. Toroid sheath approach is dead for this $G$; do not
proceed.''  This is the formal answer.

\paragraph{Why this answer is not quite a verdict on the physics.}

The $f_1$-zero outcome here reflects two things: (i) $|F|$ sits at or
below the centroid-quadrature noise floor of our BEM
($\mathcal O(h^2)\sim 10^{-4}$, varying by factors of $\sim 5$ between
$n_t=10$ and $n_t=25$); and (ii) the bootstrap on the polynomial
residuals is itself unstable -- the bootstrap-mean of $f_1$ disagrees
with the point estimate by factors $3$--$10$, so the $3\sigma$
threshold itself is not credible without a better kernel.  In other
words, our BEM at this kernel order does not resolve $f_1$ either way,
and the formal ``no signal'' conclusion should be read as
\emph{inconclusive at this resolution} rather than as a positive
statement that $f_1=0$.

\paragraph{Hypothetical: what the data would say if $f_1$ were
resolvable.}

If we set the noise-floor concern aside and treated the unconstrained
fits at face value, the $\ell$ data would unambiguously route to the
\emph{second} row of the decision table: $\ell_0\approx 17$--$21$ at
each $\alpha$ vs.\ $\ell_1\beta\lesssim 0.15$ across the swept range,
with $\mathrm{RSS}_{\ell_0=0}/\mathrm{RSS}_{\rm free}\sim 10^5$--$10^7$.
This is row 2, ``Inductance floor kills the gain.''  Note that this is
a conditional statement: it only applies if and when a future,
higher-accuracy kernel resolves $f_1>3\sigma$.

\paragraph{Caveat on the ``floor''.}

The large $\ell_0$ is \emph{entirely} the external-lead self-inductance
of the spec's wire path (out to $s=10R$, down to $z=-2R$, return along
$s\!=\!0$): XC3 returns $L_{\rm wire,self}=18.903$ of
$L_p^{\rm total}=19.030$, i.e.\ $99.3\%$.  In physical units, with all
wire distances scaling as $R_t$, this contributes
$L_{\rm wire,self}^{\rm SI}\sim\mu_0 R_t\times 19$.  The spec explicitly
notes that the readout-wire geometry is part of $G$ and ``the answer
depends on it.''  A different $G$ in which the wire return hugs the
sheath rather than going out to $s\!=\!10R$ would not have this
$\ell_0$, and could plausibly route to the first row (favorable
scaling).  The conclusion below is therefore qualified to the spec's
chosen $G$.

\section{Crossover radius $R_t^\ast$}
\label{sec:cross}

Even taking the unconstrained (potentially noise-driven) $f_1$ values
optimistically and substituting into the favorable-scaling crossover
formula gives $Q^{\rm tor}/Q^{\rm core}<10^{-12}$ at every cost budget
$V_{\rm mag}\in\{0.1,1,10\}\,\mathrm m^3$ and every $\alpha$:

\begin{table}[h]
\centering
\begin{tabular}{lllll}\toprule
$\alpha$ & $V_{\rm mag}$\,[m$^3$] & $R_t^\ast$\,[m] & $Q^{\rm tor}/Q^{\rm core}$ \\\midrule
$0.5$ & $0.1$  & $0.56$ & $5.0\!\times\!10^{-18}$ \\
$0.5$ & $1.0$  & $1.20$ & $5.0\!\times\!10^{-18}$ \\
$0.5$ & $10.0$ & $2.59$ & $5.0\!\times\!10^{-18}$ \\
$1.0$ & $0.1$  & $0.44$ & $2.0\!\times\!10^{-13}$ \\
$1.0$ & $1.0$  & $0.96$ & $2.0\!\times\!10^{-13}$ \\
$1.0$ & $10.0$ & $2.06$ & $2.0\!\times\!10^{-13}$ \\
$2.0$ & $0.1$  & $0.35$ & $5.7\!\times\!10^{-14}$ \\
$2.0$ & $1.0$  & $0.76$ & $5.7\!\times\!10^{-14}$ \\
$2.0$ & $10.0$ & $1.63$ & $5.7\!\times\!10^{-14}$ \\\bottomrule
\end{tabular}
\caption{Crossover analysis.  $Q^{\rm tor}/Q^{\rm core}\ll 1$ at every
$\alpha$ and budget because the wire-floor $\ell_0$ kills the gain.
The numerical ratio depends quartically on the unresolved $f_1$ and is
not robust as a value, but its sub-unity order of magnitude is robust:
even pretending the entire sweep-range $|F|$ were a clean signal would
not flip the inequality.}
\label{tab:crossover}
\end{table}

\section{Decision}
\label{sec:decision}

\paragraph{Formal answer per the spec's decision table.}

For all three $\alpha\in\{0.5,1.0,2.0\}$: \textbf{row 3, ``No
signal-coupled response.  Toroid sheath approach is dead for this $G$;
do not proceed''} -- because $|f_1|<3\sigma_{\rm boot}$ on every fit.
This is the formal output of the procedure laid out in
\texttt{cc\_task\_spec.md} Sec.~1.

\paragraph{Honest physical reading.}

The procedure routes to row~3 because $V_{\rm eff}$ sits at the
$\mathcal O(h^2)\sim 10^{-4}$ noise floor of our centroid-quadrature
BEM kernel, so $f_1$ is unresolved either way.  The robust observable
is $\ell(\beta;\alpha)$, which is converged and shows
$\ell_0\approx 17$--$21$ -- entirely from the external readout wire's
self-inductance ($99.3\%$ of $L_p$ from XC3).  If a future,
higher-accuracy kernel resolved $f_1$ above noise, the decision would
shift to row~2 (``inductance floor'') because $\ell_0$ is many orders
of magnitude away from being consistent with zero.  Either way, the
spec's row~1 (favorable thin-annulus scaling) is not the outcome for
this $G$.

\paragraph{What ``this $G$'' means.}

The spec instructs that the readout-wire path runs out to $s=10R$,
down to $z=-2R$, and returns along the $z$-axis.  This wire path
length is $\mathcal O(20 R)$ and dominates $L_p$ at every $\beta$
in the swept range.  The spec acknowledges (Sec.~``External return
path (precise)'') that the wire is ``part of $G$ and the answer
depends on it.''  A wire path whose self-inductance scales as
$\mu_0 R_t\beta$ rather than $\mu_0 R_t$ (e.g.\ a coaxial return that
hugs the inner sheath cylinder) would not produce this $\ell_0$
floor, and a re-run with that $G$ could plausibly fall in row~1.
\emph{This is the natural next step if the toroid topology is of
ongoing interest.}

\paragraph{Inconclusive-flags raised (spec Sec.~``What counts as
inconclusive'').}  Three flags are triggered in our data:
\begin{itemize}
\item $F(\beta)$ is non-monotone and has sign changes at every
$\alpha$, violating flag 1.
\item $\chi^2/\mathrm{dof}$ for the unconstrained $F$ fit is not
statistically meaningful with $f_1$ at noise; the polynomial ansatz
cannot be verified or refuted from these data, partial flag 2.
\item No condition-number violations ($\max\mathrm{cond}=1.8\times 10^7
<10^8$), so flag 3 does not fire.
\end{itemize}
The raw points and fits are preserved in
\texttt{outputs/production\_sweep.csv} and
\texttt{outputs/fit\_results.json}.  We do not paper over the noise:
the report (\texttt{outputs/report.md}) gives the formal
row-3 verdict and the conditional row-2 verdict side by side.

\appendix

\section{Adversarial review log}
\label{app:review}

This appendix records adversarial reviews of intermediate artifacts and
the responses they triggered.

\subsection*{Review of the initial BEM solver (\texttt{bem\_solver.py}, commit \texttt{dceadea})}

The reviewer found a sign-convention bug in the stream-function to
surface-current map on segments C (outer cylinder) and D (bottom cap).
With $t$ increasing $A\!\to\!B\!\to\!C\!\to\!D$ around the meridional
rectangle and a consistent outward-from-$V_{\rm mag}$ normal
($\hat n = -\hat s, +\hat z, +\hat s, -\hat z$), the meridional tangent
in cylindrical components is $+\hat z, +\hat s, -\hat z, -\hat s$.  For
$K\cdot\hat t=(1/s)\partial_\chi\psi$ to be globally consistent across
all four segments (so that current is conserved across corner seams),
the K-coefficient signs on segments C and D must be inverted from those
on A and B.  The reviewer also flagged the choice of straight-wire
self-inductance formula vs.\ the spec's circular-loop variant; the
former is appropriate for a polyline and the constant offset does not
affect any $\beta$-dependence.

\emph{Response.}  The sign convention was corrected to use
$K=\nabla_S\psi\times\hat n$ uniformly on all four segments
(commit \texttt{e70bac2}).  After the fix:
\begin{itemize}
\item $F$ at $\alpha=\beta=1$, $w=0.03$ converged to $\approx-1.8\times
10^{-4}$ across $n_t\in\{6,12,18\}$, replacing the pre-fix oscillation
between $5\times 10^{-5}$ and $5\times 10^{-4}$.
\item XC2 ratio improved from $\sim 0.4$ to $\sim 0.02$, comfortably
passing the $<0.1$ threshold.
\item $\ell$ values stabilised to $\approx 19.03$ across mesh
resolutions.
\end{itemize}
The straight-wire formula was kept, with the note that this is the
physically correct choice for a polyline filament.

\subsection*{Review of the production-sweep $F$ data}

The fits for all three $\alpha$ values gave $|f_1|<3\sigma_{\rm boot}$.
A higher-resolution self-check at $\alpha=1$,
$\beta\in\{0.05,0.20\}$, $n_t\in\{10,15,20,25\}$ showed that $F$ does
not converge with mesh refinement at this $G$, while $\ell$ remains
stable to $0.02\%$.  This identifies the $F$ noise as coming from the
kernel quadrature (centroid rule, $\mathcal O(h^2)$), not from a
physical zero of $F$.

\subsection*{Review of an earlier draft of Sec.~\ref{sec:decision}}

A second adversarial reviewer flagged a logical error in the earlier
draft: it pivoted from the spec's row~3 (``no signal'') to row~2
(``inductance floor'') on the strength of $\ell_0$ alone, but the
spec's decision table is exclusive on the $f_1$ axis -- $f_1$
statistically zero routes unconditionally to row~3 regardless of
$\ell_0$.  The reviewer also pointed out that the bootstrap-mean of
$f_1$ disagrees with the point fit by factors $3$--$10$
(\texttt{outputs/fit\_results.json}: at $\alpha\!=\!2.0$
$\hat f_1=4.88\times 10^{-4}$ vs.\ bootstrap mean $5.06\times 10^{-5}$;
at $\alpha\!=\!0.5$ $\hat f_1=8.26\times 10^{-6}$ vs.\ bootstrap mean
$1.87\times 10^{-5}$), so the $3\sigma$ threshold itself is unreliable
until the kernel is upgraded.  Finally, the reviewer noted that
$\ell_0$ is not a fundamental ``floor'' but the spec's specific
external-wire path -- $99.3\%$ of $L_p$ at $\alpha=\beta=1$ comes from
$L_{\rm wire,self}$ (XC3 decomposition).  In physical units the wire
length scales with $R_t$, so the contribution is
$\sim\mu_0 R_t\times 19$, part of the $L_p=\mu_0 R_t \ell$
parameterization but with $\ell$ dominated by a $\beta$-independent
constant -- exactly the spec's ``ell\_0 not consistent with zero''
condition.

\emph{Response.}  Sec.~\ref{sec:decision} was rewritten to (i) give
the formal row-3 verdict cleanly, (ii) state the conditional row-2
verdict as a hypothetical that requires future kernel work to confirm,
(iii) document the wire-path caveat prominently, and (iv) list the
``inconclusive'' flags raised per spec Sec.~``What counts as
inconclusive.''  \texttt{outputs/report.md} was updated to match.

\end{document}
